\documentclass[10pt,a4paper]{amsart}
\usepackage[margin=19mm]{geometry}
\usepackage{amsmath,amssymb,mathtools}
\usepackage[scr=rsfs,cal=euler]{mathalfa}
\usepackage{xfrac}
\usepackage[unicode,hidelinks,bookmarks=false]{hyperref}
\newcommand{\ie}{i.e.\ }
\newcommand{\cf}{cf.\ }
\newcommand{\resp}{respectively\ }
\newcommand{\Subsection}{\subsection}
\providecommand{\nobreakdash}{\nobreak\hbox{-}}
\setlength{\emergencystretch}{2em}
\multlinegap0pt
\usepackage{amssymb}
\usepackage{enumerate}
\usepackage[shortlabels]{enumitem}
\usepackage{bm}
\usepackage[all]{xy}
\usepackage{tikz}
\newtheorem{proposition}{Proposition}
\newtheorem{lemma}{Lemma}
\def\A{\op{\mathbb{A}}}
\def\Gm{\op{\mathbb{G}_m}}
\def\Z{\op{\mathbb{Z}}}
\newcommand{\op}[1]{\operatorname{#1}}
\title{Kuznetsov Categories for Gauged Linear Sigma Models}
\author{David Favero, Daniel Kaplan, Tyler L. Kelly}
\date{Source: arXiv:2512.18402v1 (CC BY 4.0)}
\begin{document}
\maketitle

\noindent\emph{Source and licence.} Kuznetsov Categories for Gauged Linear Sigma Models. Version arXiv:2512.18402v1 (CC BY 4.0), distributed by arXiv under the Creative Commons Attribution 4.0 International licence, http://creativecommons.org/licenses/by/4.0/, as recorded in the arXiv licence metadata for this exact version and shown on the abstract page https://arxiv.org/abs/2512.18402v1. Full licence notice: \texttt{licenses/arxiv-2512.18402-CC-BY-4.0.txt}.\par\smallskip

\noindent\textit{Editorial scope.} Counted unit: the article's lemma w=0 (source0.tex lines 631--633). The article's complete original proof of it is delivered with it (source0.tex lines 635--642, one \texttt{proof} environment of 8 lines). No mathematics is written, repaired or re-derived: the statement and the proof are the article's own lines, in the article's own notation, and no retained line is substituted or re-flowed.  Retained uncounted as the source-local context the proof needs: lines 612--614.  Nothing was omitted from any retained span, so the package carries no omission record.  No retained line contains a citation, so the wrapper carries no bibliography and no bibliography entry.  No retained line contains a figure environment, an image-inclusion command or drawing code, so no image is delivered and the package declares no asset dependency.  Editorial formatting only, in addition: an AMS \texttt{amsart} preamble that restates the article's own declarations and macros used by the retained lines, namely the environments \texttt{proposition}, \texttt{lemma} on separate counters, together with the standard packages those lines need.  The retained environments print in this document as Proposition 1, Lemma 1 in order of appearance, whereas the article prints the same items with the numbering of its own full document.  The complete native source of the article as distributed in the arXiv e-print for this exact version is bundled unchanged as \texttt{sources/191388/source0.tex}.
\begin{proposition}\label{prop:strongly convex weight cone}
Suppose that the divisors $D_i$ are ample for all $i$. For any wall $W \in \Sigma^{\mathcal E}_{\op{GKZ}}$, the cone $C^{\mathcal E}_W$ is strongly convex.
\end{proposition}

\begin{lemma} \label{lem: w=0}
If the divisors $D_i$ are ample for all $i$, then the function $w|_{(V \times \mathbb C^r)^\lambda}=0$. 
\end{lemma}

\begin{proof}  The function $w \in S = \op{Sym} (V \times \mathbb C^r)^\vee$ has degree $(0,1)\in \widehat{ G \times \Gm}$.  Its restriction to $(V \times \mathbb C^r)^\lambda$ is an element of degree $(0,1) \in \widehat{G/\lambda \times \Gm}$ in $\op{Sym}  ((V \times \mathbb C^r)^\lambda)^\vee$.  The polynomial ring $\op{Sym} ((V \times \mathbb C^r)^\lambda)^\vee$ has coordinates with degrees lying in $\widehat{G/\lambda} \cap C^\mathcal E_W$.  Thus, since $C^\mathcal E_W$ is strongly convex by Proposition~\ref{prop:strongly convex weight cone}, $w|_{(V \times \mathbb C^r)^\lambda}$ is constant.  However, $w|_{(V \times \mathbb C^r)^\lambda}$ has non-trivial degree in $\widehat \Gm$ and is therefore $0$.
% Take any monomial $m$ in $w$. It can be written in the form
% $$
% u_i \prod_{\rho \in R} x_\rho^{m_{i\rho}}
% $$
% for some $c \in \Gm$, $i \in \{1, \dots, r\}$, $R \subseteq \Sigma(1)$, and $m_{i\rho} \in \Z_{>0}$. If $\pi(e_i)$ or $\pi(e_\rho)$ for any $\rho \in R$ is not in the linear span $V_W$ of the wall $W$, then the monomial vanishes in $w_W$. 
% Since $w$ is homogeneous of degree $(0,1)$ in $S_{\Sigma} \times \Gm$, it is also homogeneous of degree $(0,1)$ in $S_W \times \Gm$. Thus $m$ is as well. However, the weights form a strongly convex cone $C_W^{\mathcal{E}}$ in $V_W$, by Proposition~\ref{prop:strongly convex weight cone}, hence either $\pi(e_i) \notin V_W$ or  there must be a $\rho \in R$ with $\pi(e_\rho) \notin V_W$, hence $m|_{\A^{J_W}} = 0$. This proves the claim.
\end{proof}

\end{document}
